% =====================================================================
%  Paul V. Rubow, GEORG BRANDES OG UNIVERSITETSFILOSOFIEN.
%  Danske Studier 1928 (Kjøbenhavn: Gyldendal), pp. 145--162.
%
%  TRANSCRIPTION -- GENERATED; DO NOT EDIT BY HAND.  Rebuild from the repo root:
%    python3 ebook-method/brandes-universitetsfilosofien/vote.py            (the check; must leave nothing unsettled)
%    python3 ebook-method/brandes-universitetsfilosofien/fromreads.py --write
%  Scan: ~/bibliotek/Rubow, Paul/1928.pdf (SCANS.tsv), the whole 1928 volume of Danske Studier, 300 ppi 1-bit, ABBYY
%  text layer.  printed p = PDF page - 3 (p. 145 = PDF 148, p. 162 = PDF 165).
%
%  METHOD AND ACCURACY (2026-10-09).  No e-book and no earlier transcription.  Unlike the other books in this
%  collection, the text is NOT built from the scan's text layer: it is a BLIND transcription of every page made at the
%  image by two readers (one for pp. 145-153, one for pp. 154-162) who saw no other text, checked word by word against
%  TWO independent machine readings of the same pages, the ABBYY layer and tesseract (dan) (vote.py).  A reader's word
%  that agreed with neither machine reading (25 words, 39 marks, mostly footnote marks the machines glue to words) was
%  looked at in a crop and settled; 7 corrections to the readers' text came of it (vote-fixes.tsv).  A word that both
%  machines read and the readers lack was searched for as well: none, beyond the running heads and sheet signatures,
%  which are not transcribed.  So every word rests on two of three independent readings or on a look at the image.
%  Footnotes, their marks, italics, paragraphs and the small-type quotations are the readers' (the machines cannot
%  place them).  The draft built from the layer (draft.py) and its collation against tesseract (decisions.tsv) were
%  kept as a cross-check and are not the source of this file.
%  What this cannot see: a reading in which a reader and one machine share the same fault (for instance a misprint
%  both silently corrected).
%  PRINTED AS IS: „Afhængighed at“ (p. 151) for af; „har spiller“ (p. 152) for han.  „maskeret-|teologiske“ (p. 160, in
%  a quotation from Brandes) is set with the hyphen, as a compound.  Notes are numbered by page, as printed.
% =====================================================================
\documentclass[12pt,a4paper,openany]{book}
\usepackage[utf8]{inputenc}
\usepackage[T1]{fontenc}
\usepackage{libertinus}
\usepackage{libertinust1math}
\usepackage{textalpha}
\usepackage{graphicx}
\DeclareUnicodeCharacter{0254}{\reflectbox{c}}  % ɔ (U+0254) = reversed c
\usepackage[danish]{babel}
\usepackage[protrusion=true,expansion=true]{microtype}
\usepackage{setspace}
\usepackage[margin=1.3in]{geometry}
\usepackage{fancyhdr}
\usepackage{hyperref}
\hypersetup{pdftitle={Georg Brandes og Universitetsfilosofien}, pdfauthor={Paul V. Rubow}, hidelinks}
\setlength{\headheight}{15pt}
\pagestyle{fancy}
\fancyhf{}
\fancyhead[LE,RO]{\thepage}
\fancyhead[C]{\textit{P. V. Rubow: Georg Brandes og Universitetsfilosofien (1928)}}
\setstretch{1.3}
\usepackage{marginnote}
\newcommand{\opage}[1]{\marginnote{\footnotesize\textit{[#1]}}}

\begin{document}

% ===== TEXT (Danske Studier 1928, pp. 145--162) =====
% --- p. 145 ---
\opage{145}\setcounter{footnote}{0}%

\begin{center}GEORG BRANDES OG UNIVERSITETSFILOSOFIEN\end{center}

\begin{center}AF\end{center}

\begin{center}PAUL V. RUBOW\end{center}

\begin{center}\rule{3em}{0.4pt}\end{center}

I en lille Bog, \textit{Georg Brandes og hans Lærere} (Kbh. 1927), har
jeg for nylig søgt at give et Rids af hans Udviklingsgang, eller
jeg maa snarere sige af hans „Variationer“, ved at lade den store
Kritiker passere forbi en Række mer eller mindre betydelige litterære Personer, som til forskellige Tider har paavirket ham. Efter
en Metode som sund Fornuft og Sainte-Beuves altid lærerige Exempel indskærper, har jeg dèr skelnet mellem Ungdomslærernes
Indflydelse og de senere Tilskyndelser. Det er af overordentlig
Vigtighed, naar man behandler en Forfatter, at holde ham fast i
den Periode hans Personlighed bliver til, og intet Øjeblik at slippe
ham af Syne, mens man for de senere Stadiers Vedkommende,
efter at hans Aand er bleven mere stationær, kan tillade sig at
fæste Opmærksomheden navnlig ved visse Hvile- og Vendepunkter.

Mens jeg i det omtalte Skrift, hvor Alt kom an paa den klare
Oversigt, for nogle Afsnits Vedkommende var saa heldig at kunne
henvise til Præmisser og Dokumentation i et tidligere Arbejde (se
nu mine \textit{Litterære Studier,} Kbh. 1928, S. 38 ff.), er der andre
Stykker, hvor jeg i højere Grad har maattet forlade mig paa Læserens gode Tro og Dømmekraft. Det gælder ganske særlig de første Blade, om Georg Brandes’ Forhold til den hjemlige Universitetsfilosofi, hvor det vel egentlig er forvovent at vente Kendskab
til det store litterære Efterladenskab efter de to mest indflydelsesrige Universitetsprofessorer, Rasmus Nielsen og Hans Brøchner. I
det Følgende skal jeg forsøge at give en kort Karakteristik af disse
Tænkere og Forfattere --- rigtig nok ikke nogen ganske uinteresseret Skildring, thi den er bestandig gjort med Georg Brandes’
Skrifter in mente, og med størst Opmærksomhed for den Side de
vender ud imod ham.

% --- p. 146 ---
\opage{146}\setcounter{footnote}{0}%

\begin{center}I.\end{center}

\emph{Rasmus Nielsen} var, den Gang Brandes kom til Universitetet, i sin fuldeste Kraft. Han havde netop da befriet sig for sin
tidligere næsten slaviske Hængen ved Søren Kierkegaards Tanker,
og samtidig med at hans Religionsfilosofi begyndte at antage en
ejendommelig Form, søgte han at reformere Metafysikken ved at
drage Nytte af de exakte Videnskaber.

Den Art af Tænkning han især dyrkede, var Konkordatet.
Han vilde forene de mest uensartede Foreteelser; og Formlen for
Foreningen var, at \emph{netop fordi} f. Ex. (og først og fremmest)
Tro og Viden er absolut uensartede, kan de trives i samme Bevidsthed. Var hans Dialektik ikke dybsindig, saa har den været
problemskærpende i høj Grad. I Steden for at udjevne og mediere
søgte han altid nye Uoverensstemmelser og overdrev Modsætningerne indtil det absurde. Han traadte herved i sin ældste Mesters
Fodspor, da jo Hegels Metode er anlagt efter Modsætningsprincippet; men han har just sin Force, hvor Hegel har sin Svaghed,
idet denne lægger mere Vægt paa at vise, hvorledes den ene Kategori antydes i den anden, end paa at den ene er selvmodsigende
uden den anden kommer til.

Men han løste kun Problemerne ved at hugge Knuden over.
Tro og Viden er absolut uensartede, til en Begyndelse. Dernæst er
Vilje og Erkendelsesevne, deres Organer, absolut uensartede. Saa
skelnes til Sidst mellem teoretisk og praktisk Erkendelse, og selv
her konstateres samme Uensartethed! Rasmus Nielsen var under
Indflydelse fra de forskelligste Sider: fra Martensen gik han over
til Søren Kierkegaard, og Søren Kierkegaards og Grundtvigs Tanker søgte han senere at bringe under samme Hat. Men ved Siden
af denne kristne Spekulation trivedes hos ham en radikal naturvidenskabelig Interesse.

Sansen for Fagvidenskaberne var noget overraskende i dansk
Filosofi, saa nylig efter at Martensens og Kierkegaards naturløse Tænkning ganske havde optaget Gemytterne. Rasmus
Nielsen ansaa den hegelske Spekulation for mislykket, fordi dens
Begreber ikke var tilstrækkelig kontrollerede af de enkelte videnskabelige Dicipliner. Han ønskede en ny Videnskabernes Encyklopædi til at afløse den hegelske, som endnu syntes ham belastet
med Teologi og Skolastik. Fra Slutningen af Halvtredserne arbejdede han, der hidtil kun havde tænkt over religiøse Spørgsmaal,
% --- p. 147 ---
\opage{147}\setcounter{footnote}{0}%
efter Evne med Matematik, Fysik, Kemi og Biologi. Efter nogle
mindre heldige Tilløb\footnote{\textit{Philosophisk Propædeutik i Grundtræk} (1857), \textit{Philosophie og Mathematik} (s. A.), \textit{Mathematik og Dialektik} (1859), \textit{Om Begrebet Nøiagtighed} (1860), \textit{Hr. Professor Steens „Afregning“} (s. A.).} kom han omsider godt i Gang, og hans
Serie \textit{Forelæsninger over philosophisk Propædeutik} (tre Bind fra
Begyndelsen af Treserne), \textit{Naturphilosophie} (1873) og \textit{Almindelig
Videnskabslære} (1880) viser en bestandig intimere Indarbejdelse af
Fagvidenskabernes Principper og Resultater i Filosofien.

Georg Brandes har ikke i sine unge Dage interesseret sig stort
for Naturforskning og exakte Videnskaber, og han kunde neppe
blive fængslet af Rasmus Nielsen paa Grund af hans Fortrin i
denne Retning. Men det har uden Tvivl virket vækkende paa ham
at se alle Videnskaberne bidrage til Uddannelsen af et filosofisk
System, modsat den tvetydige „immanente“ Logik hos Hegel og
Brøchner. Og senere, da han mødte den positivistiske Filosofi paa
sin Vej, var han i visse Maader forberedt.

Ved sin Alsidighed og Evne til at opstille Modsætninger fik
Rasmus Nielsen stor Betydning for mange som Indfører i Filosofiens Problemer. Georg Brandes omtaler ham selv som Vækker
med Anerkendelse og Ros. Han havde desuden et aandfuldt skriftligt og navnlig mundtligt Foredrag: man vilde vide, at han var Danmarks mest veltalende Mand.

Han indtog Brandes ved sin Personlighed. Han var i Besiddelse af en Aandslivlighed, som synes at have mindet om Brandes' egen. Han var i bestandig Udvikling; og var han ikke meget
original, saa optog han efter Haanden en Række betydelige Aanders Tanker paa selvstændig Vis. Der var i ham en evig Appetit
paa alle aandelige Foreteelser, fra religiøse til matematiske Problemer; det nye Testamentes Mirakler, Darwinismen og Forsvarssagen satte ham i lige stærke Svingninger. For den som havde
Sans for det Karakteristiske, synes han ogsaa i sin ydre Apparition at have frembudt noget Interessant, som det kan læses i
Brandes’ Nekrolog.

I en anden Henseende maa Georg Brandes have følt sig tiltrukken til ham som til en beslægtet Natur. Rasmus Nielsen var
en betydelig literær Begavelse, endda med et vist journalistisk
Anlæg. Han havde fra sin første Fremtræden viist et ubedrageligt
Blik for, hvad der var aktuelt og kunde gøre sig. Hans Filosofi er
% --- p. 148 ---
\opage{148}\setcounter{footnote}{0}%
bleven til ved Optagelse i rette Tid af Andres Tanker. Teorien
om Forholdet mellem Tro og Viden er „frit efter“ Søren Kierkegaard. Har denne Bearbejdelse ikke megen filosofisk, saa har den
til Gengæld nogen litterær Værdi. Slagordet om den absolute
Uensartethed af Tro og Viden var et udmærket Program for en
aandelig Bevægelse i 1860’erne, en glimrende Popularisation af
Kierkegaards Ideer, som i deres ægte Skikkelse ikke netop var
skikkede til at slaa an hos Borgerskabet. Og det var en dygtig
politisk Bedrift at nærme Kierkegaards og Grundtvigs Tilhængere
til hinanden.\footnote{Smlgn. Frederik Jungersen: \textit{Dansk Protestantisme ved S. Kjerkegård, N. F. S. Grundtvig og R. Nielsen,} Kbh. 1873.} Det var ogsaa nyttigt at arbejde paa at fjerne den
religiøse Mistanke fra den naturvidenskabelige Forskning. Ved sin
mangeaarige Kamp for at rive Teologien ned og paa den anden
Side bevare den protestantiske Kristendom og Kirke er han mærkværdig, snarere som Folkefører end som kritisk Tænker.

Vil man se, i hvor høj Grad han var Advokat i sin Tænkning,
skal man læse nogle af hans temmelig talrige Stridsskrifter. Her
skifter hans religionsfilosofiske Standpunkt bestandig efter hans
Udvikling og efter de forskellige Modstandere han har over for
sig. Ingen Steds viser den seje Bondenatur sig bedre end i hans
haardnakkede og ikke ganske ufrugtbare filosofiske Polemik. Hvilken Tone der skal anslaaes, hvor vidtløftig der skal gaas til Værks,
det afhænger alt af den Modspiller han har at gøre med. Biskop
Martensen er han altid parat til at maale Skæppen fuld med indtil halvandet hundrede Sider; Provst Zeuthen faar strax mindre;
Brøchner, hans Kollega i allersnevreste Forstand, hvis meget bitre
Angreb ikke kan lades ubesvaret, expederes paa godt tre Ark;
Georg Brandes og A. C. Larsen, der kun er unge Kandidater, opnaar slet intet Svar.

Han er ingen kritisk Begavelse; den indgaaende logiske Detaildiskussion af Bevisførelsen gaar han uden om, og holder sig
til visse Punkter hos Modstanderen, som kan give ham selv Opinionen paa sin Side: hos Martensen gøres det teologiske Universalmonarki til Latter, hos Brøchner den uforstaaelige Terminologi;
paa Martensen kastes den Mistanke, at han vil den frie Forskning
til Livs, Brøchner beskyldes for personligt Fjendskab og Had til
Kristendommen. Der er intet i Vejen for, at han kan gaa uden
% --- p. 149 ---
\opage{149}\setcounter{footnote}{0}%
om Hovedsagen, f. Ex. forbigaas „den absolute Uensartethed“ i
Skriftet mod Brøchner.

Rasmus Nielsen var endelig et Stykke af en Stilist. Hans
Skrivemaade er lige saa lidt oprindelig som hans Tænkning. Den
stammer direkte ned fra den i Kierkegaards filosofiske Skrifter:
den samme Overlæsselse af Klimaxer, Gentagelser, Paralleler, Antiteser, Ordspil, den hele barbariske Pynt. Men han naaede efterhaanden til at omdanne Mesterens lidenskabelige og individuelle
Foredrag til en bred, anskuelig Overtalelseskunst. Tidt misbruger
han Modsætnings-Figuren indtil Ækelhed; og han driver et Spil
med den anaforiske Sætning, som der maa stærke Nerver til ikke at
irriteres over. I \textit{Om den gode Villie} lader han paa to Sider hvert
af de fem Afsnit begynde med Spørgsmaalet: „Men hvordan skal
Dualismen hæves?“\footnote{til Dels let varieret, se anf. Skr. S. 73 ff.}. Det kan ikke nægtes, at hans Udtryksform
vinder i Pædagogik, hvad den taber i Smag. Hvad han rører ved,
faar Liv. Han forstaar at fremstille de mest abstruse Metaphysica
paa inciterende Maade og med en Slags Fattelighed, for saa vidt
som det dunkelt tænkte nogen Sinde kan blive det klart sagte.
Han er heri ganske forskellig fra Brøchner, der kvæler alt Liv i
sit logiske terminologiske Panser, og hvis indviklede Sætningsbygning ikke synes at have andet Princip end at adskille det for Opfattelsen naturligt sammenhørende.

Rasmus Nielsen skyer intet Middel til at holde Opmærksomheden fangen: folkelige Talemaader, djærve Allusioner, kendte Citater, drøje Billeder i Mængde. Selv Kunstordene indlader han sig
paa at forny; og naar han bruger de afskrækkende fremmede Gloser, personificerer han dem ofte en lille Smule, hvilket unegteligt
er gavnligere for den anskuende Opfattelse end for den logiske
Proces. Han holder i det Hele taget af at gøre Begreber til Personer eller haandgribelige Ting: „Den i sit Paradis saa ungdomsfriske, tillidsfulde, bibelstolte Skrifttheologie, hvor er den efter
Syndefaldet ved det Logisk-Physiske efterhaanden bleven svækket,
bøiet, afbleget, ydmyget!“\footnote{Anf. Skr. S. 48 f.}. Han ynder Udraab, Dialogismer og
rene Dialoger, ikke blot af A og B og Cajus, men ogsaa f. Ex.
mellem Teologien og Menigheden:

\medskip

\begin{quote}\small

„Hvad har Du“, spørger Menigheden den ved sit Syndefald syndige
Theologie, „gjort af mit Juleevangelium?“ Ak, kjære Menighed, bliv ikke
% --- p. 150 ---
\opage{150}\setcounter{footnote}{0}%
vred, Kritiken har desværre revet Juleevangeliet reent istykker. Men tab
ikke Modet! husk paa, at det er Viden, der skal hjælpe paa Troen;
husk paa, hvad Biskop Martensen siger med Anselmus: „Ikke vover jeg,
o Herre, at udgrunde dine Dybheder;“ men jeg troer, at Viden skal
hjælpe paa Troen. Græd ikke; Det Evangelium, Kritiken har revet istykker, kan Apologetiken jo sætte sammen igjen . . . .

Du seer saa mørk paa mig. Hvad kan jeg gjøre for, at Videnskaben er objectiv? Hvem skulde nu ogsaa have tænkt, at de historiske
Kjendsgjerninger i Juleevangeliet var saa forunderligt skjøre, de naturlige Bestanddele ligesaa skjøre som de overnaturlige . . . . \footnote{Ib. S. 49.}

\medskip

\end{quote}

Naar ikke Gejsten alt for meget løber af med ham, forstaar
han ogsaa at bygge en Bog simpelt og almenfatteligt op; det er
især Tilfældet med hans populære Skrifter efter 1870, men ogsaa
med tidligere Arbejder. Det gælder f. Ex. allerede \textit{En undersøgende
Anmeldelse} (1849) og hans kraftigste Stridsskrift, det før citerede
\textit{Om den gode Villie} (1867), som simpelt og klart godtgør Umuligheden af Teologiens Supremati over de andre Videnskaber ved at
paavise først Historiens, saa Naturfagenes og til sidst Filosofiens
uholdbare Stilling i en saadan Union. Ogsaa Dobbeltværket \textit{Religionsphilosophie} (1869) --- \textit{Naturphilosophie} (1873) er kløgtig anlagt,
\textit{Grundideernes Logik} (1864---66) derimod forvirret.

Overensstemmende hermed ses Georg Brandes da ogsaa i sine
smaa Anmeldelser fra Slutningen af Treserne at fornøje sig over
Rasmus Nielsen som Forfatter. I Anmeldelsen af \textit{Om den gode
Villie} tillader han sig smaa Bemærkninger til Skriftets litterære
Karakteristik: „en bred og drøj, en haanende og bidende lille Bog,
skreven med overstrømmende Frembringelseslyst og fuld af gode
Indfald. Dens Lystighed har dog paa sine Steder en mere hollandsk
end dansk Karakter“\footnote{Georg Brandes, \textit{Samlede Skrifter} XIII S. 85.}, „det udmærket vittige lille Stykke, der
ender med de Ord . . . .“\footnote{ib. 87.}. Hele Brandes’ Fremstilling af Rasmus
Nielsens Forhold til Martensen og Teologien hører samme Steds hen.

Da Brandes’ tidligere Modstander, Rasmus Nielsens frafaldne
Discipel, P. S. V. Heegaard skrev en Afhandling imod den fordums
forgudede Mesters Lære, havde Brandes den Triumf at kunne
sætte Heegaard i Rette for Mangel paa Anerkendelse af „hans
saa højt begavede Lærer“ og indflette nogle rosende Karakteristika:
„han har den yppigt frugtbare Aand, det geniale Blik, den
sjældne Oprindelighed, der er ligesaa væsentlige Evner hos en
% --- p. 151 ---
\opage{151}\setcounter{footnote}{0}%
Tænker, som Tankens strengt videnskabelige Dannelse og den
skarpe paalidelige Forstand“\footnote{ib. S. 108.}.

Det er næppe uforsvarligt at finde nogen Afhængighed at  % PRINTED AS IS: „Afhængighed at“ for „af“ (layer and tesseract read „at“; seen in a crop)
Rasmus Nielsens Stil i disse Anmeldelser og i Pjecen \textit{Dualismen i
vor nyeste Philosophi.} Man naar ikke ved det blotte Studium af
Taine og ved Gabriel Sibberns Assistance\footnote{Se mine \textit{Litterære Studier} (Kbh. 1928) S. 45 f.} til at skrive en moderne og almenfattelig Prosa. I det mindste smager Georg Brandes’
Tilbøjelighed til dialogisk Fremstilling i disse Arbejder lige saa
meget af Rasmus Nielsen som af Taine: „[det var] god Politik af
den filosofiske Kæmper at gaa angrebsvis til Værks og besvare
Biskop Martensens høflige „Goddag“ med et frygteligt „Økseskaft“\footnote{G. B. \textit{Saml. Skr.} XIII S. 91.}.
„Se til \textit{Grundideernes Logik,} siger Forfatteren [ɔ: Heegaard], det
Bedste var, om Professor Nielsen, der dog er ude af Stand til
at bygge den færdig, straks lod den staa hen som en „interessant
Ruin“, altsaa som en Slags Marmorkirke. „Selv Marmorkirke“, svarer Professor Nielsen, rolig pegende paa første Del af den vidtberømte \textit{Indledning til den rationelle Etik}“\footnote{ib. S. 107.} (i dette Skrift var Heegaard som bekendt endnu paa Nielsen’ske Stadier, og det blev
vitterlig aldrig fortsat). Her findes Vittigheder i utvivlsom Nielsen’sk
Manér: „Professor Nielsen har vel én Mand i Vagten, men han græder --- med det ene Øje og ler med det andet; hans højre Haand
véd ikke, hvad den venstre gør. Han bevæger sig frem paa to
uensartede Ben, et Kødben og et Træben, paa to absolut spaltede
Ben, flækket som den ene af Uffes Saksere. Hans Hoved er et Janushoved: med den ene Mund forkynder han Naturlovens højere
Opfyldelse igennem Miraklet; med sin anden Strube giver han
Kulturbevidstheden Æren og synger en velment Lovsang for David
Strauss“\footnote{ib. S. 91 --- i selve Sagens Realitet er der ikke noget overdrevet: de som i vore Dage mener, at Jesus aldrig har eksisteret som historisk og faktisk Person, kan uden Risiko sætte Rasmus Nielsen paa Listen over deres Meningsfæller.}.

Hvor dybt Indtrykket af Rasmus Nielsen stikker, ses til sidst
tydelig af den Nekrolog Brandes skrev over ham\footnote{\textit{Saml. Skr.} II S. 419 ff.} (1884). Her
giver han ham ganske sikkert som Personlighed alt hvad han
fortjente, som Tænker lidt til, og fremhæver fint og sikkert hans
oversete Fortrin som Skribent.

Som vi har set, stod han altsaa i nogen Taknemlighedsgæld
til sin gamle Filosofikumsprofessor. Han regnedes af Samtidige fra
% --- p. 152 ---
\opage{152}\setcounter{footnote}{0}%
alle Lejre med megen Grund for at høre til Rasmus Nielsens farligste Fjender. Det skaffede ham Artigheder af Nielsens ældre Modstandere og paadrog ham Forbandelser af Tilhængerne. Saaledes
nævnes han med Honnør i Martensens \textit{Om Tro og Viden} (1867)
--- hvis Diktion ligesom næsten alle Indlæg i denne berømmelige Fejde
er en Smule smittet af Rasmus Nielsens --- som den „begavede
Forfatter“, hvem man „kun maa ønske, at han ikke alt for tidligt
maa komme ud over det Platoniske \textit{θαυμαζεῖν} (Forundringen) som
Philosophiens Begyndelse“\footnote{Anf. Skr. S. 8.}. Omvendt faar Brandes baade læst og
paaskrevet i den førnævnte Bog af Fr. Jungersen, \textit{Dansk Protestantisme,} hvor Pjècen om \textit{Dualismen} kaldes „en overordentlig
letfærdig Kritik“\footnote{Anf. Skr. S. 210 f.}. De mange Gange Brandes selv i Skrifter eller
i Interviews omtalte Rasmus Nielsen, var det altid med Højagtelse,
til Tider næsten med Ømhed. Studerer man Forholdet paa nært
Hold, ser man at der fører mere end een Traad fra den store
voltairianske Kritiker tilbage til den kierkegaardianske Sofist.

\begin{center}II.\end{center}

Forholdet til \emph{Hans Brøchner} er langt simplere og gennemsigtigere. Den Mand Georg Brandes tilegnede sit Hovedværk, var
hans egentlige strenge Opdrager; har spiller en lignende Rolle i  % PRINTED AS IS: „har spiller“ for „han spiller“ (seen in a crop, .render/*-odd-01.png)
hans Liv som Vacherot i Taines og var som Vacherot ikke nogen
ukritisk Tilskuer til Elevens hastige Udvikling.\footnote{I \textit{Tilskueren} for Februar 1912 har Dr. Poul Levin S. 213 meddelt Brøchners Censur over Brandes' indsendte Afhandling for Doktorgraden, som jeg her tillader mig at aftrykke: „Det forekommer mig, at Cand. Brandes's Afhandling, naar den bedømmes fra et rent philosophisk Synspunkt, har to ikke uvæsentlige Mangler. Behandlingen af de Problemer, der drages frem, er noget aphoristisk, og i Paavisningen af Taine's Forhold til sine historiske Forudsætninger paa Philosophiens Gebet er der en Ufuldstændighed, der gjør, at det ikke bliver tilstrækkeligt iøjnefaldende, i hvor paafaldende en Grad han som Philosoph mangler Originalitet, Dybde og Consequents. Men disse Mangler ved Afhandlingen opveies vistnok tilfulde ved dens store Fortrin: ved dens rige Indhold, der vidner om en omfattende og en vel tilegnet Kundskabsfylde hos Forfatteren, ved den skarpe, fine og aandfuldt individualiserede Opfattelse af Gjenstanden, og ved Fremstillingens Liv og Friskhed. Jeg nærer derfor ikke nogen Tvivl om, at Afhandlingen bør antages.“} Brandes har i sit
\textit{Levned} med Fromhed noteret Dagen for sit første Besøg hos den
elskede Lærer: „det var den 2. September 1861“ og i alle Enkeltheder gjort Rede for deres intellektuelle Samliv. Her skal kun tilføjes hvad en Kommentator yderligere kunde have at erindre.

Brøchner var en stilfærdig, indadvendt Tænker og Lærd, som
% --- p. 153 ---
\opage{153}\setcounter{footnote}{0}%
den Gang endnu ikke havde udgivet sine vigtigste Værker. Men
han havde dem færdige i Hovedet og kunde senere publicere dem i
endegyldig Form i Løbet af forbavsende kort Tid. Gennem Forelæsninger og privat Samtale var deres Indhold Brandes bekendt
længe før de kom i Trykken.

I et Privatbrev til Chr. K. F. Molbech har Brøchner fremsat det
som sit Ønske at efterlade to Arbejder: et der kunde vise, han havde
vidst noget, et der kunde vidne om at han havde forstaaet at tænke.

Det første har han lagt for Dagen i sine Arbejder til Filosofiens Historie, hvor han dels --- med en gammel Generalstabsofficers Sikkerhed --- skeletterer de mest repræsentative metafysiske
Bygningsværker fra den græske Filosofi i dens tragiske Tidsalder
indtil Hegel, dels spekulativt læser Verdensaandens Arbejde i de
Love han mener bestemmer Tænkningens Udvikling, og som Tænkerne ubevidst retter sig efter. En Litteraturhistoriker har begribeligvis meget at udsætte paa disse Lærdomsværker paa Grund af
deres stivnakkede deduktive Fremgangsmaade --- Brøchner havde
en Forkærlighed for Skrifter i Form af Haandbøger og altomfattende Systemer, og en mageløs Foragt for Sammenhængen mellem
Filosofien og de øvrige Videnskaber, som gør Partier af hans Skrifter,
f. Ex. Behandlingen af den franske Materialisme i det 18. Aarhundrede, aldeles misvisende; og hans Tilbøjelighed til at etablere
spekulative Sammenhænge svækker Betydningen af de Afsnit hvor
han var bedst, ogsaa specielt filologisk, funderet, som i Fremstillingen af Grækernes Filosofi.\footnote{Hans Indsigt i Middelalderens Tænkning er kun kommet Andre til Gode gennem Kommentaren til Chr. K. F. Molbechs Danteoversættelse (1851---62), hvortil han har leveret det filosofiske Materiale.} Men hans historiske Arbejder,
navnlig Haandbogen i to Bind (1873---74) aftvinger endnu Respekt
ved deres Præcision og Grundighed.

Sin Evne til selvstændig Tænkning har han godtgjort i de to
Bøger, som giver hans egen Filosofi: \textit{Problemet om Tro og Viden}
(1868) og \textit{Om det Religiøse i dets Enhed med det Humane} (1869).

Han begynder --- ikke uden Sans for det Aktuelle --- med en
Kritik af den samtidige danske Filosofis og Teologis Forsøg paa at
løse Spørgsmaalet Tro contra Viden. Han afviser den teologiske
Trosfilosofi, repræsenteret af Martensen, som vil opbygge en overlogisk Videnskab om Ting, han selv indrømmer ikke kan vides.
Han afviser fremdeles med Ærbødighed og Beundring Søren Kierkegaard, der bryder Sammenhængen af mellem Personligheden og
% --- p. 154 ---
\opage{154}\setcounter{footnote}{0}%
den almengyldige Erkendelse ved at bestemme Troen som „den
objektive Uvished, fastholdt i den mest lidenskabelige Inderligheds
Tilegnelse“, og som ved at fornegte Menneskets Naturside sætter
uopløselig Disharmoni i Personligheden og Tilværelsen. Han afviser endelig Rasmus Nielsens Lære om den absolute Uensartethed:
som en Skinløsning, et Anarki af uforenelige Modsigelser.

Vil man forklare Forholdet mellem Tro og Viden, maa man
ifølge Brøchner tage sit Udgangspunkt i en af disse Sfærer. Og
han vælger --- som Hegel --- Viden, fordi kun den kan give en
selvstændig Opfattelse af Tilværelsen. Troen kunde ikke staa alene,
men maatte støtte sig til Viden over for en saadan Opgave.

Han bygger nu en hel monistisk og idealistisk Metafysik op.
Han forkaster een Gang for alle den atomistiske Synsmaade og
vil kun operere med Totaliteter. Hele Tilværelsen er en absolut
Totalitet, og inden for den er atter Individet en relativ Totalitet,
som har sin Opgave at løse inden for den absolute.

Tilværelsens Princip er ideelt. Det Reale er et Gennemgangsstadium i Ideens Udvikling; Ideen er baade Totalitetens Grund og
dens Endemaal. Brøchner opløser alle Dualismer: Erkendelse og
Vilje er ikke uensartede, Viljen har en Erkendelse til Indhold, og
den udvikler sig „fra den første umiddelbare, enkelte Driftbestemthed“ gennem Reflexion „til en Bestemmethed ved bevidst dannede
Grundsætninger, Maximer, af hvilke saa en Livsanskuelses Enhed
kan fremgaae“\footnote{\textit{Om det Religiøse} S. 18 f.}. Heller ikke er der nogen Kvalitetsforskel mellem
teoretisk og praktisk Erkendelse. Man sætter Forskellen i, at den
praktiske Erkendelse er personlig, den teoretiske upersonlig --- men
„Erkjendelsesacten, der er underkastet Sandhedsideens Norm, har i
denne et ‘Bør’, i Forhold til hvilket der er Bevidsthed om Ansvar“\footnote{ib. S. 15.}.

Brøchners Religionsfilosofi er bestemt ved disse Grundbegreber.
Det religiøse Forhold ses som Menneskeaandens „absolute Forhold
til det Absolute“ ɔ: dens udelte og ubetingede Hengivelse til og
Subordination under den evige og uforanderlige Grundkraft i Tilværelsen. Baade Erkendelse og Vilje maa være med. Erkendelsen
fortsættes i et Postulat: Troen. Der findes i al lavere og højere
Erkendelse et Trosmoment, nemlig Tillid til Enheden og Lovmæssigheden i Tilværelsen, og et Hengivelsesmoment: „I al sand Erkjendelse er der en Tankens Luttring, en Opgiven af individuel
Hildethed, Menen og Fordom, med eet Ord: af Tankens Selvisk%
% --- p. 155 ---
\opage{155}\setcounter{footnote}{0}%
hed; der er i den en Selvhengivelse til Sandhedens absolute Magt“\footnote{ib. S. 50 f.}.
Den Tro som er til Stede i Erkendelsen, er rationel, og maa derfor forkaste alt hvad der strider mod Begrundelsens og Nødvendighedens Princip, f. Ex. Miraklet. Viljen fortsættes i Selvhengivelse,
idet Individet bliver sig bevidst „som organisk Led i Tilværelsens
Totalitet“ og „vælger at handle derefter“. Trosmomentet viser sig
i det Etiske, som er „Udtrykket for den sande, væsentlige Villen
og Handlen“; det kommer frem „i enhver ethisk Beslutning, i
hvilken Individet bestemmer sig til Handling, trods Erkjendelsens
Ufuldstændighed og Virkeliggjørelsens, Udfaldets Usikkerhed“\footnote{ib. S. 55.}.
Hengivelsesmomentet ses i, at det Etiskes Princip er Kærlighed.

Den Strid der bestandig raser mellem Repræsentanter for Viden og for Troen i dens positive Form, skyldes da ikke principiel
Modsætning: Aarsagen maa ligge i Menneskeaandens Endelighedsvilkaar. Ser vi bort fra alt det Endelige i de positive Religioner,
finder vi, at det vigtigste eller ene Fornødne i dem er Kravet om
Subordination, Selvhengivelse, som altsaa ogsaa findes i Viden.
Tro og Viden strider da ikke mod hinanden: der bliver „vel ligesom et Mellemrum mellem Tro og Viden, men et Mellemrum, der
maa kunne gjennemløbes af Viden“\footnote{ib. S. 51.}. Brøchner mente i Religionernes Historie at kunne iagttage et bestandigt Fremskridt i Retning af mere lutret Religiøsitet. Det højeste Standpunkt er Humanitetens Religion, som ikke behøver at skabes, men blot at paavises,
thi det Religiøse findes allerede, „om end ofte uklart og halvbevidst“ i vor Tids humane Bevidsthed.

I Forbindelse med sin religiøse Opfattelse behandler Brøchner
de etiske Problemer i meget spekulativ Aand. Den frie Vilje ses i
Menneskets Evne til at hævde sig efter sin uendelige Væsensbestemthed over for det Ydre og Tilfældige. Det Onde er en Forstyrrelse i „det normale Forhold mellem Menneskets Væren som
relativ Totalitet og som Led i den universelle Totalitet, og mellem
Natur- og Aandssiden i det.“ Viljen er ond, naar den er Egenvilje
og sætter sig op imod det Absolutes Fordring. Paafaldende er
denne etiske Opfattelses Optimisme: der lades vidtstrakte Muligheder aabne for Anger og Forsoning og udtales Mistillid til den
onde Viljes Bæreevne. Det Onde er nærmest en Tilfældighed, som
tilmed har størst Sandsynlighed for at korrigere sig selv. Viljen til
det Onde har en naturlig Tilbøjelighed til at slaa om til det Gode;
% --- p. 156 ---
\opage{156}\setcounter{footnote}{0}%
den laaner sin Kraft fra det Absolute\footnote{Sammenlign Georg Brandes’ \textit{Æsthetiske Studier} S. 116: „det Slettes Magt er den Kraft, det har suget af det Gode.“}, og maa saaledes kæmpe
mod sit eget Væsens Grund og til Slut ødelægge sig selv.

Til Syvende og Sidst er hans Anskuelse et religiøst Lyssyn.
Værdierne bevares; intet Væsentligt gaar tabt. Han nægter f. Ex.
personlig Udødelighed, men mener at den individuelle Aand ved
Organismens Tilintetgørelse optages „som Moment i det Eviges
Liv“; den forvandles til „et Moment i Guds Viden af sig i sit
evige Aandsrige“\footnote{\textit{Om det Religiøse} S. 189.}. Tilværelsen er Kosmos.

Brøchner bebrejdede den hegelske Filosofi, at den ikke lod
Virkeligheden komme tilstrækkelig til sin Ret. Det samme kan
imidlertid siges om hans eget System. Det er spiritualistisk og
transscenderende som Hegels. Det er uden Forhold til Fagvidenskaberne, for saa vidt langt mindre moderne end Rasmus Nielsens.
Det er med al sin Konsekvens og Fasthed yderst exklusivt, aldeles
afsluttet til alle Sider. Det er ikke godt at vide hvad Fordel Filosofien skulde kunne drage i Fremtiden af Naturforskning og exakt
Videnskab, skønt Brøchner selv anviste Fremtidens Filosoffer denne
Arbejdsmark i sine filosofihistoriske Skrifter. Ogsaa Metoden er den
gamle „immanente“ Begrebsudvikling, saa meget Brøchner end selv
har ivret imod den\footnote{Se hans \textit{Bidrag til Opfattelsen af Philosophiens historiske Udvikling} (1869) S. 265 f}.

Ovenstaaende tarvelige Redegørelse for Brøchners religiøse
Fritænkning --- et sidste Skud paa den romantiske Filosofis Stamme,
en principfast Epigons spekulative Testamente --- kan tjene til
Baggrund for og til Fuldstændiggørelse af den interessante Skildring
i Brandes’ \textit{Levned} af hvorledes han blev Fritænker, og kan sammenstilles med den summariske Fremstilling i min førnævnte Afhandling om \textit{Georg Brandes og hans Lærere} S. 6 ff. Det siger sig
selv, at den intime Paavirkning fra Mester til Lærling i Tiden før
Brøchner begyndte sit egentlige Kerne-Forfatterskab, og før Brandes
overhovedet gav noget i Trykken, er en Ting som ikke kan
gøres til Genstand for objektiv Undersøgelse og saglig Drøftelse.

Men man ser i det mindste af Brandes’ egne Ord og af Brøchners trykte Skrifter, hvad det var for en Slags Opdragelse Naturalismens tilkommende Kritiker modtog. Hans ungdommelige
Tænkning er bleven skolet i en abstrakt og idealistisk Begrebsfilosofi. Og de Tænkere som Brøchner indførte ham i, var af samme
% --- p. 157 ---
\opage{157}\setcounter{footnote}{0}%
Art. Men uanset alt hvad Brandes senere lod falde som i Grunden
uforeneligt med sit Væsen, er der Træk i Brøchners Filosofi, som
maatte vinde hans fulde Tilslutning. Den er, som viist, strengt
monistisk og modvirkede alle Fristelser som Tidsalderen havde nok
af, til en dualistisk Livsanskuelse, ved sin ubetingede Fastholden
af eet Princip. Dernæst er den udpræget intellektualistisk. Den hævder Erkendelsens Principat i Menneskeaanden, og vælger Viden til
at give Nøglen til Tilværelsens Gaader. Endelig er den erklæret
fritænkerisk. Brøchner var en haardnakket Forkæmper for „den
frie Tankes“ ubegrænsede Rettigheder og blev lagt for Had som
Fjende af Kristendommen, en Beskyldning han med megen Kraft
og Værdighed afviser i sit \textit{Svar til Professor R. Nielsen} (1868).
Hans Mening om Kristendommen indses af hans Syn paa positive
Religioner i det Hele taget; desuden indeholder hans Bog \textit{Om det
Religiøse} rundt om skarp Kritik af de kristelige Hovedforestillinger:
Miraklet, Skabelsen, Udødeligheden, og af den kristnes Læres Tendens til Fornægtelse af det Naturlige --- han holdt paa Rigtigheden
af Kierkegaards Kristendomsbillede imod Martensens. Han havde
som ung maattet døje en Afvisning fra teologisk Embedsexamen
for sin aabent vedkendte Opfattelse af Miraklerne som mytiske.
Han var en af dem der havde indført Strauss\footnote{\textit{Den christelige Troeslære \dots{} Oversat af H. Brøchner, Cand. phil} 2 Bd. 1842---43.} og gransket Feuerbach, en Tænker som har betydet meget for Udviklingen af det religiøse Standpunkt han urokkelig holdt fast ved.

De Skrifter af Georg Brandes, hvori man nærmest vil søge
specielle og haandgribelige Paavirkninger fra Brøchner, er \textit{Den
franske Æstetik i vore Dage,} hans mest akademiske Arbejde, og
\textit{Emigrantlitteraturen,} hans Fritænker-Manifest. Men det er mikroskopisk hvad her er at finde.

Disputatsen, hvori han saa at sige fremlægger sin yngre Lærers Doktriner for sin ældre Lærer, er visselig et renligt Elevarbejde, som følger Mesterens Metode til at præparere et filosofisk
System, give dets Forudsætninger og fælde Dommen (hvorom de
to Parter, som det vil ses af Noten ovf. S. 152, dog altsaa ikke
var enige). Men i det Enkelte finder man kun Lapperier. Det vigtigste Sted er S. 241, hvor jeg finder Paavirkning fra Brøchners
\textit{Om det Religiøse:} „Taine er gaaet ud fra Spinozas Synspunkt, at
Mennesket er en Del af det Hele og ikke ‘en Stat i Staten’; men
% --- p. 158 ---
\opage{158}\setcounter{footnote}{0}%
han er gaaet saa vidt, at han har sluttet, at Mennesket heller ikke
kan være Borger i denne Stat, kun er et Produkt uden Personlighed, uden Frihed og uden virkelig Ret. Han kommer saaledes
til at ignorere Selvbevidstheden af den gode Grund, at han ikke
kan forklare den, idet han mangler ethvert Begreb om \textit{Mennesket
som relativ Totalitet} og forstaar Naturen ud fra Mechanismens
istedenfor fra Organismens Idee.“

Ogsaa i de tre første Bind af Hovedstrømningerne har Brandes
ønsket at fortsætte Brøchners Gerning; man var klar over det i
Samtiden, og Modstanderne haabede at ramme Mesteren, naar de
skød efter Discipelen. \textit{Emigrantlitteraturen}\footnote{Om hvis akademiske Karakter og hele Stilling jeg kan henvise til mine \textit{Litterære Studier} S. 80 ff.} er tilegnet Brøchner
„med Taknemmelighed og Venskab“. \textit{Den romantiske Skole i Tyskland} har følgende, ikke senere optrykte, Efterskrift:

\medskip

\begin{quote}\small

Kjære Hr. Professor!

\medskip

Det er en Selvfølge, at dette Bind lige saa fuldt som det foregaaende
er tilegnet Dem. Det indeholder med mine Universitetsforelæsninger som
Grundlag Adskilligt, som ikke kunde finde Plads eller ikke vilde være paa
sin Plads i det mundtlige Foredrag. Det har været mit Haab under Udarbeidelsen, at Mænd, som De, vilde finde Behag i denne Bog. En Bog
er for mig en Handling. De vil bedømme den som saadan.

De kjender Heines Ord, Hr. Professor: „Sonderbar! und immer ist
es die Religion, und immer die Moral, und immer der Patriotismus, womit alle schlechten Subjecte ihre Angriffe beschönigen! sie greifen uns
an, nicht aus schäbigen Privatinteressen, nicht aus Schriftstellerneid, nicht
aus angeborenem Knechtsinn, sondern um den lieben Gott, um die guten
Sitten und das Vaterland zu retten.“ --- Lykkeligvis opveies for mig en
hel Pakke af anonyme Artikler og Brochurer af et Ord fra Dem.
\par\hfill Kjøbenhavn, April 1873.\par\hfill Deres\par\hfill G. B.

\medskip

\end{quote}

Men skønt \textit{Hovedstrømningerne} efter sit oprindelige Anlæg er
en Fortsættelse af Brøchners Polemik med Teologer og Dualister ---
selve Titelen og det ofte tilbagevendende Slagord „den frie Tanke“
stammer fra ham --- saa tror jeg ikke det vil være muligt i det
Enkelte at paavise brøchnerske Paavirkninger deri. Det er først i
sjette Bind, Bogen om \textit{det unge Tyskland,} at den allerede aldrende
Brandes knytter Traaden til sin hegelianske Ungdom og skildrer
den unghegelianske Kamp for Tankens Frihed og imod Klerikal%
% --- p. 159 ---
\opage{159}\setcounter{footnote}{0}%
ismen. Han indsætter her sin Lærers venstrehegelianske Fortolkning
af Hegels Metafysik\footnote{Sammenlign atter \textit{Georg Brandes og hans Lærere} S. 13 ff.}.

Der er kun en ganske lille Gruppe af Brandes’ Skrifter, hvor
man kan sige Paavirkningen fra Brøchner gaar mere i det Specielle:
det er de polemiske Arbejder fra Tiden paa begge Sider af 1870.
For Resten er det muligt, at Paavirkningen er gensidig; i al Fald
staar Brøchners ældste Lejlighedsskrift, et Indlæg i Daabs-Fejden
i Begyndelsen af Fyrrerne meget tilbage for dem, hvormed han
bidrog til Striden om Tro og Viden. Det var Georg Brandes selv,
der tvang ham Pennen i Haand\footnote{Se Brandes’ \textit{Levned} I S. 177.} og fik ham til først at kritisere
de Samtidiges Opfattelser, dernæst udvikle sin egen; saa gik Resten
af sig selv.

I \textit{Dualismen i vor nyeste Philosophi} (1866), som er ældre end
Brøchners Indlæg i Fejden, skal man følgelig ikke vente formel
Paavirkning, men Tankegangen er ganske sikkert helt igennem inspireret af Brøchners, den stemmer ogsaa fortræffeligt med Brøchners
egne Deduktioner to Aar senere i \textit{Problemet om Tro og Viden.} Først
gives der en Fremstilling af alt det Fristende i Rasmus Nielsens
Filosofi: der kastes endog Mistanke paa den eventuelle Modstander:
„den Enkelte, der optræder ikke for Sagens Skyld, men for at
bringe sit eget lille Jeg i Forgrunden ved Hjælp af den store Sag“\footnote{Anf. Skr. S. 11.}.
Men saa „vendes Medaljen“, og de explosive Stoffer nedkastes med
saa meget større Kraft. Og efter at Tilliden til Systemet er rystet,
bevises med streng Logik, baade at Præmisserne til den opstillede
Thesis ikke stemmer med Konklusionen, og at denne i og for sig
er en Selvmodsigelse. Dernæst paavises det, at Hjælpebegrebet
„det Antirationelle“ er en haabløst obskurantistisk Forestilling.
Endelig demonstreres den paastaaede Dualismes ødelæggende Følger
i hele Ophavsmandens System, og de mislykkede Sammenlodningsforsøg blottes, indtil Systemet til sidst fordømmes i sin Ynkelighed
som et Stykke mislykket Eklekticisme.

Grundtankerne er følgende: 1\textsuperscript{o} --- Trossætninger og Videnssætninger, som Rasmus Nielsen opstiller dem, er ikke absolut uensartede. Sætninger som „Verden er skabt i sex Dage“ --- „Verden
er ikke skabt i sex Dage“ er blot stridige, uforenelige. Men selv
om Tro og Viden virkelig var uensartede, hjalp det dog Intet, da
Uensartethed kun er et andet Udtryk for Uforenelighed. Og det
% --- p. 160 ---
\opage{160}\setcounter{footnote}{0}%
vilde forudsætte en Spaltning baade i Bevidstheden og Tilværelsen.
2\textsuperscript{o} --- Trospostulater kan kritiseres af Viden: „den som --- det være
sig nu som troende eller vidende --- antager Nogetsomhelst om
Jordens Tilblivelse, han har en Antagelse, som Geologien kan corrigere“\footnote{ib. S. 31.}. 3\textsuperscript{o} --- Det „Antirationelle“, Rasmus Nielsens mageløse
Opdagelse, falder fuldstændig uden for videnskabelig Diskussion, da
det ikke kan udtrykkes i rationelle Begreber eller blot i Sproget.

Religionen kritiseres ikke i det lille Skrift. En Ytring lader
Muligheden af mirakelfri Religion staa aaben: „Hvis ‘Religionen
forudsætter en Aabenbaring og Aabenbaringen et Under, men Videnskaben fornegter Underet og forkaster Aabenbaringen’ . . . saa
kan den Form af det Religiøse, der lader sig forene med Videnskaben, kun være en saadan, der er uden Undere og uden Aabenbaring“\footnote{ib. S. 60. Jfr. Stedet om Parker ib. S. 73 f.}.

\textit{Forklaring og Forsvar, En Antikritik} (1872) er omvendt i
Form, men neppe i Indhold, paavirket af de mellemliggende polemiske Skrifter af Brøchner, navnlig af \textit{Et Svar til Professor
R. Nielsen.} Brandes beundrede i højeste Grad Brøchners Duelighed
som Polemiker. I sin Afhandling om den elskede Lærer fra 1886
skriver han herom:

\medskip

\begin{quote}\small

Hans Angrebskrig mod Teologerne eller mod Nielsens maskeret-teologiske Filosofi var et Mønster paa overlegen Taktik. Han udviklede
gerne først med største Nøjagtighed den Anskuelse, han vilde bestride,
fremsatte dernæst Punkt for Punkt Beviserne for dens Uholdbarhed, anførte saa Alt, hvad der fra Modstanderens Side kunde tænkes gjort gældende mod disse Beviser, og forklarede, hvorfor dette formodede Forsvar
Intet vilde betyde; saa beviste han, at andet Forsvar end det anførte
overhovedet ikke kunde føres.

Naar dette var gjort, opløste han Modstanderens Begrebssammenhæng,
gjorde ham Indrømmelser, satte \textit{per impossibile} at et af de kuldkastede
Hovedpunkter var et fast Kastel, og viste, at saa ligefuldt de mest afgørende Begrebsbestemmelser indeholdt Modsigelse paa Modsigelse. Skridt
for Skridt gik han rundt om sin Modstanders Lærebygning, og vel bekendt
med dennes Kampmaade stoppede han ethvert Hul, hvorigennem han
kunde ville undslippe, indtil han havde ham sikkert fangen. Saa endte
han gerne med at opregne Hovedpunkterne i alle de paaviste Modsigelser,
og sluttede tilsidst den i den alvorligste Tone holdte Kritik med en Ironi,
som man i al Sandhed og uden Misbrug af Ordet kunde kalde frygtelig.\footnote{\textit{Saml. Skr.} II S. 429 f.}

\medskip

\end{quote}

% --- p. 161 ---
\opage{161}\setcounter{footnote}{0}%
Brandes’ egen Taktik i de fem smaa Kapitler af hans \textit{Forklaring}
af 1872 er ikke meget forskellig fra den han her har beskrevet.
Især er det tydeligt, at han efterligner Brøchner hver Gang han
fra Defensiven gaar over til Offensiven. Han imødegaar først de
talrige Presseangreb. Ved behændigt valgte Citater kommer Modstanderne til at slaa sig selv for Munden, saa ofte de fremsætter
en Paastand; og han er ikke bange for at benytte sin Sejr saadan,
at han giver Modstanderen Ret i, hvad der nylig er bevist at være
Uret, for saa at godtgøre, at det alligevel ikke hjælper ham noget:

\medskip

\begin{quote}\small

Jeg har sagt: Det, at en Litteratur i vore Dage lever, viser sig i,
at den sætter Problemer under Debat. Jeg har, hvad \textit{Dagbladet} uden
videre forbigaar, sagt: i vore Dage, jeg har ikke paastaaet, at paa Homers
eller Hesiods Tid laa Beviset for Litteraturens Liv i dens Discussion af
Problemer. Jeg har lige saa lidt talt om Shakespeares Tidsalder. Men
naar min Anmelder i \textit{Dagbladet} endelig spørger, hvilke Problemer \textit{Lear}
eller \textit{Macbeth} sætter under Debat, da tillader jeg mig uden videre at henvise ham, som en Begynder i det æsthetiske Studium, til den første den
bedste Haandbog om Shakespeare, f. Ex. Ulricis’s eller Gervinus’s. Der
vil han finde Svar. . . .\footnote{Anf. Skr. S. 20.}

\medskip

\end{quote}

Krigen føres over i Fjendens eget Land. Rudolph Schmidt har
bebrejdet Georg Brandes hans Efterligninger i Indhold og Sprog
af fransk Litteratur:

\begin{quote}\small

Hvad Ondt er der saa deri? Naar man gaar i Solen, bliver man
solbrændt. Men sæt nu hine Anklager kom fra en Litterat, der har fulgt
sin Lærer gjennem Tykt og Tyndt som en Lakai følger sin Herre, hvis
Mening er den Mening Hr. N. har idag, hvis Mening var den, som Hr. N.
havde igaar, og hvis Mening vil være den, som Hr. N. vil have imorgen!
Han glemmer da at sige, at medens jeg idetmindste har indaandet Europas frie Luft, har han kun indaandet sin Lærers Aande.\footnote{ib. S. 39 f. --- R. Schmidt skrev under Mærket „k“; Brandes bryder, som man ser, Pseudonymiteten, og alluderer til hans trofaste Væbnertjeneste hos „N.“, ɔ: Rasmus Nielsen.}

\medskip

\end{quote}

Men i Slutningen af hvert Kapitel hæver han sig fra det letvundne
Mundhuggeri til almindelige Principper. Særlig gælder dette Polemikken mod „44“, et Navneskjul for Biskop Monrad. Endelig
slutter han, ligesom Brøchner i sit Svar til Rasmus Nielsen, med
en værdig Fremhævelse af at hans Overbevisning er fast og dyrekøbt og derfor ikke let at vælte over Ende ved Angreb.

\medskip

% --- p. 162 ---
\opage{162}\setcounter{footnote}{0}%
De ovenstaaende Iagttagelser og Overvejelser er ganske vist
kun Smaating. Men naar de gælder et Forfatterskab, hvis Plads
i dansk Litteraturhistorie ganske sikkert altid vil blive meget fremstaaende, er de dog ikke overflødige. Det er altid Umagen værd
at mindes Navnene paa de ejendommelige Personligheder, som har
indført den unge Georg Brandes i Aandslivets Problemer, og at
holde dem fast under Læsningen af hans tidligste Skrifter.

\end{document}
